%% Complete-positivity regions of Runge--Kutta discretizations of phase-covariant qubit dynamics. BIT Numerical Mathematics; Springer Nature sn-jnl class.
\documentclass[pdflatex,sn-mathphys-num]{sn-jnl}
\usepackage[T1]{fontenc}
\usepackage{graphicx}
\usepackage{amsmath,amssymb,amsfonts}
\usepackage{amsthm}
\usepackage{mathtools}
\usepackage{booktabs}
\usepackage{algorithm}
\usepackage{algorithmicx}
\usepackage{algpseudocode}
\usepackage{xcolor}
\hypersetup{pdftitle={Complete-positivity regions of Runge\textendash Kutta discretizations of phase-covariant qubit dynamics},pdfauthor={G. Blake Pierpoint; Olivier Bernard; Yichen Liu},pdfsubject={Research article},pdfkeywords={Lindblad equation; complete positivity; Runge\textendash Kutta methods; stability regions; Richardson extrapolation; certified numerics}}
% PDF outline: sections and subsections only, so that run-in heads do not nest into a staircase.
\hypersetup{bookmarksdepth=2}

%% Notation macros
\newcommand{\CPTP}{\mathrm{CPTP}}
\newcommand{\CP}{\mathrm{CP}}
\newcommand{\id}{\operatorname{id}}
\newcommand{\sgn}{\operatorname{sgn}}
\newcommand{\Tr}{\operatorname{Tr}}
\newcommand{\spec}{\operatorname{spec}}
\renewcommand{\Re}{\operatorname{Re}}
\newcommand{\e}{\mathrm{e}}
\newcommand{\ii}{\mathrm{i}}
\newcommand{\Fmargin}{F_{R,\theta}}
\newcommand{\Mmargin}{M_{R,\varpi}}
\newcommand{\Adm}{\mathcal A}
\newcommand{\Lgen}{\mathcal L}
\newcommand{\Ldis}{\mathcal L_{\mathrm{dis}}}
\newcommand{\Uh}{\mathcal U}
\newcommand{\PASS}{\textnormal{\textsc{accept}}}
\newcommand{\FAIL}{\textnormal{\textsc{reject}}}
\newcommand{\UNCERTAIN}{\textnormal{\textsc{undecided}}}

\theoremstyle{thmstyleone}
\newtheorem{theorem}{Theorem}
\newtheorem{proposition}[theorem]{Proposition}
\newtheorem{corollary}[theorem]{Corollary}
\newtheorem{lemma}[theorem]{Lemma}
\theoremstyle{thmstylethree}
\newtheorem{definition}{Definition}
\theoremstyle{thmstyletwo}
\newtheorem{remark}{Remark}
\newtheorem{example}{Example}

\raggedbottom
% Break URLs only after slashes and similar characters, never after the scheme colon.
\AtBeginDocument{\def\UrlBigBreaks{}\def\UrlNoBreaks{\do\(\do\[\do\{\do\<\do\:}}

\begin{document}

\title[Complete-positivity regions of Runge--Kutta discretizations]{Complete-positivity regions of Runge--Kutta discretizations of phase-covariant qubit dynamics}

\author*[1]{\fnm{G. Blake} \sur{Pierpoint}}\email{pierpogb@odu.edu}
\author[2,3]{\fnm{Olivier} \sur{Bernard}}\email{oli.bernard@weylcenter.org}
\author[4,3]{\fnm{Yichen} \sur{Liu}}\email{Yichen.Liu@weylcenter.org}

\affil*[1]{\orgname{Old Dominion University}, \orgaddress{\street{5115 Hampton Boulevard}, \city{Norfolk}, \state{Virginia} \postcode{23529}, \country{USA}}}
\affil[2]{\orgdiv{UFR PhITEM}, \orgname{Universit\'e Grenoble Alpes}, \orgaddress{\city{Saint-Martin-d'H\`eres}, \country{France}}}
\affil[3]{\orgname{Weyl Center for Mathematical Physics}, \orgaddress{\city{Washington}, \state{DC}, \country{USA}}}
\affil[4]{\orgdiv{Department of Computer Science and Engineering}, \orgname{Shanghai Jiao Tong University}, \orgaddress{\city{Shanghai}, \country{China}}}

\abstract{A Runge--Kutta step of a Lindblad equation can be absolutely stable yet not completely positive. For the autonomous phase-covariant qubit generator, two values of the stability function fix the update, so the Choi conditions decide exactly which steps are channels for every method with real coefficients. On one-way, zero-dephasing rays the exact flow lies on the channel boundary, and a signed defect law gives the leading margin coefficient as the first exponential defect times a polynomial in the precession-to-relaxation ratio; when nonzero, its sign decides small-step admissibility, also for Richardson extrapolates. On these rays the admissible set of the classical fourth-order method has six regimes as the ratio varies, including disconnected sets and isolated steps. The rotating frame restores the zero-precession set, one full step and two half steps can have disjoint positive-step admissible sets, and on the one-way, zero-dephasing ray at zero precession the fifth-order Richardson extrapolate of the classical method is a channel at no positive step where both constituents are. Disconnectedness persists under small dephasing or bidirectional relaxation and is certified at three strengths of a transverse field that breaks phase covariance. For phase-covariant candidates an exact acceptance test certifies the specified update at the executed double-precision step and, in adaptive tests, rejects non-channel updates that error control accepts.}

\keywords{Lindblad equation, complete positivity, Runge--Kutta methods, stability regions, Richardson extrapolation, certified numerics}

\pacs[MSC Classification]{65L06, 65L20, 65L05, 65G20, 81S22, 81P47}

\maketitle

\section{Introduction}\label{sec:intro}

A practitioner who integrates a Markovian master equation with an unchanged Runge--Kutta code obtains, at every step, a linear map on density matrices. The exact flow $\e^{t\Lgen}$ of a Gorini--Kossakowski--Sudarshan--Lindblad (GKSL) generator $\Lgen$ is completely positive and trace preserving (CPTP) for every $t\ge0$ \cite{Lindblad1976,GoriniKossakowskiSudarshan1976,BreuerPetruccione2002}; the numerical map need not be. The requirement on that map is complete positivity, a property of the map and not of one trajectory: a linear map $\Phi$ on $2\times2$ matrices is completely positive if and only if its Choi matrix $J(\Phi)=\sum_{jk}|j\rangle\langle k|\otimes\Phi(|j\rangle\langle k|)$ is positive semidefinite \cite{Choi1975CP,Kraus1971}. A map that keeps every input density matrix positive can still send an entangled input to an output that is not positive semidefinite, so a single computed trajectory cannot establish complete positivity when the update is reused for many inputs or applied to one subsystem of a larger register.

Applied to the linear system $\dot\rho=\Lgen\rho$, an $s$-stage Runge--Kutta (RK) method with stability function $R$ advances the state by the map $R(h\Lgen)$. Riesch and Jirauschek \cite{RieschJirauschek2019} showed that the RK and predictor--corrector update maps are not CPTP in general, whereas the matrix exponential is. One response is to build a different update. Kraus-form, integrating-factor, nested-Picard, Gregory-quadrature, tensor-train, quantum-channel, and exponential Euler discretizations preserve structure by construction \cite{RieschPiklJirauschek2020,CaoLu2025,AppeloCheng2025,HuAppeloCheng2025NPI,HuAppeloCheng2026Gregory,DelMastroAppeloCheng2026TT,RobinRouchonSellem2025,ChenEtAl2026}, with different guarantees: preserving density matrices after a state-dependent trace normalization is not the same as defining a linear CPTP map \cite{RobinRouchonSellem2025}. Classical numerical analysis offers thresholds instead. Positivity, contractivity, absolute monotonicity, and strong-stability preservation give step-size bounds under which an RK method preserves a convex constraint on the solution vector \cite{BolleyCrouzeix1978,Kraaijevanger1991Contractivity,Horvath2005Threshold,GottliebShuTadmor2001}; exponential RK methods of order above two cannot preserve positivity throughout the linear problem class, including time-dependent inhomogeneity, considered in \cite{OstermannVanDaele2007}; and positivity can be enforced adaptively by perturbing the RK weights \cite{NussleinRanochaKetcheson2021}. None of this answers the practitioner's question: for a fixed generator and a fixed way of building the update from RK stages, which step sizes produce a quantum channel?

We answer it exactly for the autonomous phase-covariant qubit generator, the family combining coherent precession about the $z$ axis, downward and upward relaxation, and pure dephasing, whose channel geometry is classical \cite{RuskaiSzarekWerner2002,Filippov2020Phase,Siudzinska2023Geometry}. The generator acts diagonally on a population mode and a coherence mode, so the RK map is determined by two scalar values of the stability function, and the Choi conditions for phase-covariant channels pull back to explicit inequalities in the step size. Complete positivity implies absolute stability of each mode for the same linear map, but the converse fails: the channel inequality couples the modes, and the admissible set can be disconnected or contain isolated points. Threshold results of the type in \mbox{\cite{BolleyCrouzeix1978,Kraaijevanger1991Contractivity,Horvath2005Threshold}} cannot describe this, since a threshold is the maximal interval attached to zero. The amplification-matrix stability analysis of Gregory nested-Picard schemes \cite{HuAppeloCheng2026Gregory} concerns a different constructed update without trace renormalization, not the Choi admissibility of the unchanged RK map.

The criterion itself (Theorem~\ref{thm:criterion}) holds for every RK method with real coefficients. On the one-way, zero-dephasing ray the exact flow lies on the boundary of the channel set, so the discretization error decides on which side the numerical step falls. Theorem~\ref{thm:defect} makes this quantitative: the leading coefficient of the complete-positivity margin is $(-1)^m\eta_mG_m(\varpi)$, where $\eta_mz^m$ is the first defect of $R(z)-\e^z$ and $G_m$ is an explicit polynomial in the precession-to-relaxation ratio $\varpi$. When the coefficient is nonzero its sign, the small-step orientation, decides admissibility of small steps, and zeros of odd multiplicity of $G_m$ reverse the orientation as the frequency varies, so the order of a method does not by itself determine whether small steps are channels. Corollary~\ref{cor:extrapolation} extends the law to $n$-fold Richardson extrapolates of any method. The admissible set of the classical fourth-order method has six regimes as $\varpi$ varies (Theorem~\ref{thm:rk4}), among them a narrow reentrant band, a regime with no positive admissible step, two isolated admissible steps, and a detached high-frequency window whose edges we characterize asymptotically; the critical frequencies are algebraic numbers, certified by exact root isolation.

The admissible set belongs to the update, not to the method. Transforming to the rotating frame before discretizing gives a different linear map, with the zero-precession admissible set at every frequency (Proposition~\ref{prop:frame}). For the classical fourth-order method substepping dilates the admissible set, so one full step and two half steps can have disjoint positive-step sets, and in the high-frequency window, below an exact threshold frequency, repeated halving can skip the admissible component (Propositions~\ref{prop:composition}--\ref{prop:ladder}). On nonrotating one-way damping without dephasing the fifth-order Richardson extrapolate of the classical method is non-CPTP at every positive step of the interval on which both of its constituents are channels (Theorem~\ref{thm:richardson}). Nor is the disconnected structure an artefact of the boundary: it persists under small positive dephasing or bidirectional relaxation (Theorem~\ref{thm:persistence}), and exact certificates at $\theta=\kappa=10^{-3}$, $\varpi=2$ exhibit two admissible components at the transverse-field strengths $\Omega_x/\Gamma\in\{0.05,0.10,0.20\}$, where phase covariance is broken, as at the phase-covariant reference $\Omega_x=0$. The practical outcome is an acceptance test that certifies, by exact rational arithmetic, the specified update at the executed double-precision step size, treats an undecided case as a rejection, and in adaptive tests rejects the non-channel updates that error control accepts.

Sections~\ref{sec:problem}--\ref{sec:rk4} set up the problem, prove the criterion and the defect law, and classify the classical fourth-order method; Sections~\ref{sec:candidates} and~\ref{sec:guard} treat the construction of the update, persistence, and certified acceptance; Sections~\ref{sec:discussion} and~\ref{sec:conclusion} discuss scope and conclude. Our research note \cite{PierpointBernardLiu2026Note}, with code archived at \cite{PierpointBernardLiu2026Zenodo}, reported the zero-precession signed-defect law, certified admissible sets for eight RK formulas on zero-temperature amplitude damping (including $[0,\alpha_4]$ for the classical fourth-order method), the positive-dephasing buffer, and the composition and fixed-horizon observations. New here are the frequency-dependent defect law and its Richardson extension, the classification of the classical method over all frequencies, the substep ladder, the bidirectional-relaxation buffer, the persistence theorem, and the representation, transverse-field, and acceptance-test results. Certificates, sign charts, and derivations not needed to follow the proofs are in Online Resource~1, and every computation is reproducible from the code archive cited in the Declarations.
\section{Phase-covariant generator and Runge--Kutta candidates}\label{sec:problem}

\subsection{Generator, coordinates, and modes}\label{sec:generator}

Units with $\hbar=1$ are used throughout, and decimal expansions followed by an ellipsis are truncated, not rounded. Let $|0\rangle$ and $|1\rangle$ be the lower- and upper-energy states, $\sigma_-=|0\rangle\langle1|$, $\sigma_+=|1\rangle\langle0|$, and $\sigma_z=|0\rangle\langle0|-|1\rangle\langle1|$. The autonomous phase-covariant generator is
\begin{equation}
 \Lgen(\rho)=-\ii[H,\rho]+\gamma_\downarrow\mathcal D[\sigma_-](\rho)+\gamma_\uparrow\mathcal D[\sigma_+](\rho)+\frac{\gamma_\phi}{2}(\sigma_z\rho\sigma_z-\rho),
 \label{eq:gksl}
\end{equation}
with $H=-\omega\sigma_z/2$, $\mathcal D[L](\rho)=L\rho L^\dagger-\{L^\dagger L,\rho\}/2$, and nonnegative rates $\gamma_\downarrow,\gamma_\uparrow,\gamma_\phi$ \cite{Filippov2020Phase,Siudzinska2023Geometry}. Let $\Gamma=\gamma_\downarrow+\gamma_\uparrow>0$ and define the dimensionless quantities
\begin{equation}
\begin{aligned}
 \theta&=\frac{\gamma_\uparrow}{\Gamma}\in[0,1],&
 \kappa&=\frac{\gamma_\phi}{\Gamma}\ge0,&
 \varpi&=\frac{|\omega|}{\Gamma}\ge0,\\
 x&=\Gamma h,& u&=\gamma_\phi h=\kappa x,& \nu&=\omega h,
\end{aligned}
 \label{eq:coordinates}
\end{equation}
where $h>0$ is the step size; $\theta$ is the stationary excitation fraction, and $|\nu|=\varpi x$. The generator acts on the population $\rho_{11}$ and the coherence $\rho_{01}$ by
\begin{equation}
 \dot\rho_{11}=-\Gamma(\rho_{11}-\theta),\qquad
 \dot\rho_{01}=-\Big(\frac\Gamma2+\gamma_\phi-\ii\omega\Big)\rho_{01},
 \label{eq:modes}
\end{equation}
while $\Tr\rho$ is conserved. Consequently $\spec(\Lgen)=\{0,\,-\Gamma,\,-(\Gamma/2+\gamma_\phi\mp\ii\omega)\}$, and the exact semigroup multiplies the population deviation $\rho_{11}-\theta$ by $a_{\rm ex}=\e^{-x}$ and the coherence by $c_{\rm ex}=\e^{-x/2-u+\ii\nu}$. The multipliers remain well defined as $\Gamma\to0$, although the ratios $\theta$, $\kappa$, and $\varpi$ do not.

\subsection{Direct Runge--Kutta steps}\label{sec:rkstep}

For an $s$-stage RK method with real Butcher coefficients $(A,b)$ let
\begin{equation}
 R(z)=1+z\,b^{\mathsf T}(I-zA)^{-1}\mathbf 1
 \label{eq:stability_function}
\end{equation}
be its stability function \cite{HairerNorsettWanner1993}. Applied to the linear autonomous system $\dot\rho=\Lgen\rho$, one step of size $h$ is the linear map $R(h\Lgen)$, provided the stage equations are solvable, that is $\det(I-\lambda A)\ne0$ for every $\lambda\in\spec(h\Lgen)$. We call this set of steps the \emph{stage-evaluation domain}; for explicit methods it is unrestricted, and for implicit methods it excludes the poles of $R$ reached by $-x$ and $-x/2-u\pm\ii\nu$. A removable cancellation in the reduced rational expression of $R$ does not regularize a singular stage system, so all statements below are made on the stage-evaluation domain.

The methods used in the accuracy and order comparisons are consistent, $R'(0)=1$; the finite-step criterion itself does not need consistency.

Because $\Lgen$ acts diagonally on the modes \eqref{eq:modes}, the direct step is determined by two scalar values of $R$,
\begin{equation}
 a=R(-x),\qquad c=R(-x/2-u+\ii\nu),
 \label{eq:multipliers}
\end{equation}
and it acts on $X=(X_{jk})\in M_2(\mathbb C)$ as
\begin{equation}
 \Phi_\theta(a,c)(X)=
 \begin{pmatrix}
 A_\theta X_{00}+C_\theta X_{11} & cX_{01}\\
 \bar cX_{10} & B_\theta X_{00}+D_\theta X_{11}
 \end{pmatrix},
 \label{eq:direct_map}
\end{equation}
where
\begin{equation}
 A_\theta=1-\theta(1-a),\quad B_\theta=\theta(1-a),\quad
 C_\theta=(1-\theta)(1-a),\quad D_\theta=\theta+(1-\theta)a.
 \label{eq:ABCD}
\end{equation}
Since the RK coefficients are real, $a$ is real, and $\Phi_\theta(a,c)$ preserves the trace (${A_\theta+B_\theta=C_\theta+D_\theta=1}$) and Hermiticity for every step in the stage-evaluation domain. Only complete positivity is at stake. Replacing $\omega$ by $-\omega$ conjugates $c$ and leaves $|c|$ and every statement below unchanged, so only $\varpi=|\omega|/\Gamma$ enters.

\subsection{Candidate maps and representation}\label{sec:candidate_def}

An adaptive method may build several linear maps over one interval and advance only one of them. Their complete-positivity properties can differ, so we name them.

\begin{definition}[Candidate maps]\label{def:candidate}
Fix the generator \eqref{eq:gksl}, a method with stability function $R$, and a total dimensionless step $x>0$ (total dephasing step $u$, total phase $\nu$). The \emph{full-step candidate} is $\Phi_R(x)=R(h\Lgen)$, with multipliers \eqref{eq:multipliers}. The \emph{equal-substep candidate} with $n\ge2$ substeps is $\Phi_R^{(n)}(x)=\Phi_R(x/n)^n$. The \emph{extrapolated candidate} of linear order $p$ is
\begin{equation}
 E_n\Phi_R(x)=\frac{n^p\,\Phi_R^{(n)}(x)-\Phi_R(x)}{n^p-1}.
 \label{eq:extrapolated_candidate}
\end{equation}
For $n=2$ and $p=4$ this is the standard Richardson extrapolate $(16\Phi_R^{(2)}-\Phi_R)/15$. The \emph{rotating-frame candidate} is defined in Section~\ref{sec:frame}. A \emph{projected candidate} is a full-step or equal-substep candidate re-evaluated at a smaller step, and a \emph{fallback candidate} is a map produced by a different propagator over the same interval; both are specified in Section~\ref{sec:guard}.
\end{definition}

Every candidate in Definition~\ref{def:candidate} shares the block structure \eqref{eq:direct_map} with its own multipliers: composition multiplies $a$ and $c$ (Proposition~\ref{prop:composition}), and the extrapolated candidate has multipliers $E_n R(-x)$ and $E_nR(-x/2-u+\ii\nu)$ with $E_nR(z)={\big(n^pR(z/n)^n-R(z)\big)/(n^p-1)}$, because \eqref{eq:direct_map} is affine in $(a,c)$ for fixed $\theta$.

By the \emph{laboratory frame} we mean discretization of $\Lgen$ itself; by the \emph{rotating frame} we mean exact transformation of the equation to the interaction picture of $H$ followed by discretization of the transformed generator. The two produce different linear maps over the same interval; Section~\ref{sec:frame} quantifies the difference.

\subsection{Complete positivity and the admissible set}\label{sec:cp_def}

With the convention $J(\Phi)=\sum_{jk}|j\rangle\langle k|\otimes\Phi(|j\rangle\langle k|)$, the map \eqref{eq:direct_map} has the Choi matrix
\begin{equation}
 J_\theta(a,c)=
 \begin{pmatrix}
 A_\theta&0&0&c\\
 0&B_\theta&0&0\\
 0&0&C_\theta&0\\
 \bar c&0&0&D_\theta
 \end{pmatrix},
 \label{eq:choi}
\end{equation}
which is the established Choi form of a non-unital phase-covariant qubit map \mbox{\cite{Choi1975CP,RuskaiSzarekWerner2002,Siudzinska2023Geometry}}; $\Phi$ is completely positive if and only if $J\succeq0$. For a candidate family and a fixed generator ray $(\theta,\kappa,\varpi)$ we write
\begin{equation}
 \Adm_R(\varpi)=\{x\ge0:\ \Phi_R(x)\ \text{is CPTP}\},
 \label{eq:admissible_set}
\end{equation}
and $\Adm_R^{(n)}$, $\Adm_R^{\rm ext}$ for the substep and extrapolated candidates; $\Adm_4$ denotes the set of the classical fourth-order method. The step $x=0$ always belongs to these sets. We say that a set has a \emph{positive-step component} when it contains an interval of positive length, and that a component is \emph{attached} if it contains $0$ and \emph{detached} otherwise.
\section{Complete-positivity criterion and boundary orientation}\label{sec:criterion}

\subsection{The finite-step criterion}\label{sec:finite_step}

The Choi matrix \eqref{eq:choi} is small enough to diagonalize by hand, and doing so turns complete positivity into four inequalities in the two multipliers.

\begin{theorem}[Finite-step criterion]\label{thm:criterion}
Let $R$ be the stability function of a Runge--Kutta method with real coefficients, let $(x,u,\nu)$ lie in the stage-evaluation domain, and let $a=R(-x)$ and $c=R(-x/2-u+\ii\nu)$. The direct step $\Phi_\theta(a,c)$ of \eqref{eq:direct_map} is trace preserving and Hermiticity preserving, and it is completely positive if and only if
\begin{equation}
 a\le1,\qquad A_\theta\ge0,\qquad D_\theta\ge0,\qquad
 \Fmargin(x,u,\nu):=a+\theta(1-\theta)(1-a)^2-|c|^2\ge0.
 \label{eq:criterion}
\end{equation}
\end{theorem}

\begin{proof}
Trace and Hermiticity preservation were established after \eqref{eq:direct_map}. After a permutation of the basis, the Choi matrix \eqref{eq:choi} is the direct sum of the scalars $B_\theta$ and $C_\theta$ and of the Hermitian block $\big(\begin{smallmatrix}A_\theta&c\\ \bar c&D_\theta\end{smallmatrix}\big)$. For $\theta\in(0,1)$ both scalars are nonnegative exactly when $a\le1$; for $\theta=0$ only $C_\theta=1-a$ is present and for $\theta=1$ only $B_\theta=1-a$, so in every case the scalar conditions reduce to $a\le1$. The block is positive semidefinite exactly when its diagonal entries and its determinant are nonnegative, and direct multiplication gives $A_\theta D_\theta=a+\theta(1-\theta)(1-a)^2$.
\end{proof}

The criterion specializes the known Choi conditions of phase-covariant qubit channels \cite{RuskaiSzarekWerner2002,Siudzinska2023Geometry} to the stability function, and the rest of the paper studies the inequalities \eqref{eq:criterion} as functions of the step size; we call $\Fmargin$ the \emph{Choi margin}. Two consequences are used throughout. First, at the one-way endpoints $\theta\in\{0,1\}$ the criterion reduces to
\begin{equation}
 0\le a\le1,\qquad \Mmargin(x):=a-|c|^2\ge0,
 \label{eq:one_way_criterion}
\end{equation}
where the one-way margin $\Mmargin$ already implies $a\ge0$, and the lower eigenvalue of the active block is
\begin{equation}
 \lambda_-=\frac{1+a-\sqrt{(1-a)^2+4|c|^2}}{2}
 =\frac{2\Mmargin}{1+a+\sqrt{(1-a)^2+4|c|^2}}\,,
 \label{eq:lambda_margin}
\end{equation}
so its sign is the sign of $\Mmargin$; the population ceiling $a\le1$ and the block margin are independent exit mechanisms from the channel set. Complete positivity also implies absolute stability of every mode: \eqref{eq:criterion} gives $|a|\le1$ and $|c|^2\le A_\theta D_\theta\le1$, so $\Adm_R(\varpi)\subseteq\{x:\ |R(-x)|\le1,\ |R(-x/2-u+\ii\nu)|\le1\}$, and the difference between the two sets is the subject of this article. Second, the exact semigroup has $a_{\rm ex}=\e^{-x}$ and $|c_{\rm ex}|^2=\e^{-x-2u}$, hence
\begin{equation}
 F_{{\rm ex},\theta}=\e^{-x}\big(1-\e^{-2u}\big)+\theta(1-\theta)\big(1-\e^{-x}\big)^2,
 \label{eq:exact_buffer}
\end{equation}
which, for $x>0$, vanishes exactly when $u=0$ and $\theta\in\{0,1\}$. On these one-way, zero-dephasing rays, which we call the boundary rays, the exact flow lies on the boundary of the channel set, so the discretization error determines which side contains the numerical step. Elsewhere, that is when $\kappa>0$ or $\theta\in(0,1)$, the exact flow lies in the relative interior with $F_{{\rm ex},\theta}>0$, and we then speak of a \emph{buffer}. For a consistent method, $R(z)=1+z+\tau_2z^2+O(z^3)$, the Choi margin satisfies
\begin{equation}
\begin{aligned}
 \Fmargin(x,u,\nu)&=2u+O\big(x^2+xu+u^2+\nu^2\big),\\
 \Fmargin(x,0,\varpi x)&=\Big[\theta(1-\theta)+\frac{\tau_2}{2}-\frac14+(2\tau_2-1)\varpi^2\Big]x^2+O(x^3),
\end{aligned}
 \label{eq:buffers}
\end{equation}
so dephasing supplies a first-order buffer and bidirectional relaxation a second-order buffer $\theta(1-\theta)x^2$ for every method of order at least two ($\tau_2=1/2$). The dephasing-buffer mechanism was already described in our research note \cite{PierpointBernardLiu2026Note}. For every consistent method any $\kappa>0$ therefore yields a positive attached admissible interval, and for a method of order at least two so does any $\theta\in(0,1)$. Section~\ref{sec:persistence} estimates the lengths for the classical method.

\subsection{Signed first-exponential-defect law}\label{sec:defect}

We now fix the boundary rays $\theta\in\{0,1\}$, $u=0$, $\nu=\pm\varpi x$, write $\alpha=\tfrac12-\ii\varpi$, and suppose that
\begin{equation}
 R(z)-\e^z=\eta_mz^m+O(z^{m+1}),\qquad m\ge2,\qquad \eta_m\ne0,
 \label{eq:first_defect}
\end{equation}
so that $m-1$ is the linear order of the method. Define
\begin{equation}
 G_m(\varpi)=1-2\,\Re\,\alpha^m,\qquad q=1+4\varpi^2,\qquad \sigma_m=2^m(\alpha^m+\bar\alpha^m).
 \label{eq:G_definition}
\end{equation}
Because $\alpha$ and $\bar\alpha$ are the roots of $z^2-z+q/4$, the numbers $\sigma_m$ satisfy $\sigma_0=\sigma_1=2$, $\sigma_m=2\sigma_{m-1}-q\sigma_{m-2}$, and $G_m=1-\sigma_m/2^m$ is a polynomial in $q$; for instance
\begin{equation}
\begin{aligned}
 G_2&=\frac q2,& G_3&=\frac{3q}{4},& G_4&=-\frac{q(q-8)}{8},\\
 G_5&=-\frac{5q(q-4)}{16},& G_6&=\frac{q(q^2-18q+48)}{32},& G_7&=\frac{7q(q-4)^2}{64}.
\end{aligned}
 \label{eq:G_factors}
\end{equation}

\begin{theorem}[Signed first-defect law]\label{thm:defect}
Let $R$ be analytic at $0$ with real Taylor coefficients and satisfy \eqref{eq:first_defect}. Then, on the boundary rays,
\begin{equation}
 \Mmargin(x)=R(-x)-\big|R(-\alpha x)\big|^2=(-1)^m\eta_mG_m(\varpi)\,x^m+O(x^{m+1}).
 \label{eq:defect_law}
\end{equation}
Hence if $G_m(\varpi)\ne0$ there is $x_0>0$ such that for $0<x<x_0$ the direct step is CPTP when $(-1)^m\eta_mG_m(\varpi)>0$ and non-CPTP when $(-1)^m\eta_mG_m(\varpi)<0$.
\end{theorem}

\begin{proof}
From \eqref{eq:first_defect}, $R(-x)=\e^{-x}+(-1)^m\eta_mx^m+O(x^{m+1})$ and $R(-\alpha x)=\e^{-\alpha x}+(-1)^m\eta_m\alpha^mx^m+O(x^{m+1})$. Since $\Re\alpha=\tfrac12$, $|\e^{-\alpha x}|^2=\e^{-x}$, and $\e^{-\bar\alpha x}=1+O(x)$, so
$|R(-\alpha x)|^2=\e^{-x}+2(-1)^m\eta_m\,\Re(\alpha^m)\,x^m+O(x^{m+1})$. Subtracting gives \eqref{eq:defect_law}; all semigroup terms of lower order cancel exactly. For small $x$ the constraint $0\le a\le1$ holds because $a=1-x+O(x^2)$, so the sign of $\Mmargin$ decides membership by \eqref{eq:one_way_criterion}.
\end{proof}

We call the sign of the leading coefficient of $\Mmargin$ the \emph{small-step orientation} of $R$ on the ray: \emph{inward} when it is positive, so that sufficiently small steps are channels, and \emph{outward} when it is negative. The leading coefficient is $(-1)^m\eta_mG_m(\varpi)$ when $G_m(\varpi)\ne0$ and otherwise belongs to a higher power of $x$, beginning with \eqref{eq:degenerate} below. The hypothesis holds for the stability function of every consistent RK method with real coefficients, for the compositions $R(z/n)^n$, and for the extrapolates $E_nR$ of Definition~\ref{def:candidate}, so the law applies to the full-step, substep, and extrapolated candidates; the rotating-frame candidate has the margin of $\varpi=0$ (Proposition~\ref{prop:frame}). The law is the channel-inequality analogue of the local sector structure of the order star of $R$ at the origin \cite{IserlesNorsett1991OrderStars}, whose $2m$ sectors are fixed by the sign of $\Re(\eta_mz^m)$: the same first defect $\eta_m$ controls both, but here it is multiplied by the frequency factor $G_m$, whose real zeros can reverse the orientation. At $q=4$, that is $\varpi=\sqrt3/2$, $\alpha=\e^{-\ii\pi/3}$ and $G_m(\sqrt3/2)=1-2\cos(m\pi/3)$, which vanishes exactly when $m\equiv\pm1\pmod 6$; a zero of odd multiplicity reverses the orientation, whereas a zero of even multiplicity ($G_7$ at $q=4$) leaves it the same on both sides, and at the zero itself the sign is decided by the next coefficient. When $G_m(\varpi)=0$ the next coefficient is
\begin{equation}
 \big[(-1)^{m+1}\eta_{m+1}G_{m+1}(\varpi)+2(-1)^m\eta_m\,\Re(\bar\alpha\alpha^m)\big]x^{m+1}+O(x^{m+2})
 \label{eq:degenerate}
\end{equation}
(Online Resource~1, Section~2). It is consistent with the exact factorizations $M_{4,\sqrt3/2}(x)=-x^6(x-2)^2/576$ for the classical fourth-order method at $\varpi=\sqrt3/2$ and $x^5(3-2x)/9$ for the three-stage, third-order strong-stability-preserving method SSPRK(3,3) \cite{GottliebShuTadmor2001} at $\varpi=\sqrt7/2$, both obtained by direct expansion. The principal fifth-order Dormand--Prince formula \cite{DormandPrince1980} has $m=6$ and $\eta_6=1/3600$, so its small-step orientation is inward, outward, and inward again across the two zeros $0.7509\ldots$ and $1.8536\ldots$ of $G_6$, that is $\varpi={(2\mp\sqrt{33}/4)^{1/2}}$. Methods of the same order can have opposite orientations, as the two fourth-order formulas in Table~\ref{tab:methods} show. Table~\ref{tab:methods} lists the first defects and orientations at $\varpi=0$ of the methods used in this article, computed from the exact Butcher tableaux.

\begin{table}[tbp]
\caption{First exponential defect $\eta_mz^m$ and local orientation coefficient $(-1)^m\eta_mG_m(0)$ of the one-way margin at zero precession, from Theorem~\ref{thm:defect} and Corollary~\ref{cor:extrapolation}. A positive coefficient means that sufficiently small steps are CPTP. The first eight rows reproduce the zero-precession results of \cite{PierpointBernardLiu2026Note}; the Richardson row is new}\label{tab:methods}
\footnotesize
\begin{tabular*}{\textwidth}{@{\extracolsep\fill}lccccl}
\toprule
Method & $p$ & $m$ & $\eta_m$ & $(-1)^m\eta_mG_m(0)$ & Small steps \\
\midrule
Forward Euler & 1 & 2 & $-1/2$ & $-1/4$ & non-CPTP \\
Two-stage order two (Heun) & 2 & 3 & $-1/6$ & $1/8$ & CPTP \\
SSPRK(3,3) & 3 & 4 & $-1/24$ & $-7/192$ & non-CPTP \\
Classical RK4 & 4 & 5 & $-1/120$ & $1/128$ & CPTP \\
Dormand--Prince 5 (principal) & 5 & 6 & $1/3600$ & $31/115200$ & CPTP \\
Dormand--Prince 4 (embedded) & 4 & 5 & $97/120000$ & $-97/128000$ & non-CPTP \\
Backward Euler & 1 & 2 & $1/2$ & $1/4$ & CPTP \\
Implicit midpoint & 2 & 3 & $1/12$ & $-1/16$ & non-CPTP \\
Richardson extrapolate of RK4 ($n=2$) & 5 & 6 & $-1/4320$ & $-31/138240$ & non-CPTP \\
\botrule
\end{tabular*}
\end{table}

The zero-precession specialization of the signed law, $(-1)^m(1-2^{1-m})\eta_mx^m$, was stated in \cite{PierpointBernardLiu2026Note}; the frequency factor extends it to laboratory-frame precession.

\subsection{Extrapolated candidates}\label{sec:extrapolation}

Richardson extrapolation of RK steps is classical, and its effect on absolute stability, dissipation, and dispersion has been studied \cite{Szyszkowicz1985Richardson,HavasiKazemi2018Richardson}. Its effect on complete positivity follows from Theorem~\ref{thm:defect}, because extrapolation acts simply on the first defect.

\begin{corollary}[Extrapolated candidates]\label{cor:extrapolation}
Let $R$ have real coefficients and linear order $p\ge1$, $R(z)-\e^z=\eta_{p+1}z^{p+1}+\eta_{p+2}z^{p+2}+O(z^{p+3})$, and let $n\ge2$. Then:
\begin{enumerate}
 \item $E_nR(z)-\e^z=\eta^E_{p+2}z^{p+2}+O(z^{p+3})$ with
 \begin{equation}
  \eta^E_{p+2}=\frac{(n-1)(\eta_{p+1}-\eta_{p+2})}{n\,(n^p-1)};
  \label{eq:extrapolation_defect}
 \end{equation}
 in particular the extrapolated candidate has linear order $p+1$, and a higher order exactly when $\eta_{p+1}=\eta_{p+2}$.
 \item If $\eta^E_{p+2}\ne0$ and $G_{p+2}(\varpi)\ne0$, then on the boundary rays the extrapolated candidate satisfies $M_{E_nR,\varpi}(x)=(-1)^{p+2}\eta^E_{p+2}G_{p+2}(\varpi)x^{p+2}+O(x^{p+3})$, so for small steps it is CPTP if and only if $(-1)^p\eta^E_{p+2}G_{p+2}(\varpi)>0$.
 \item If $R$ is the truncated exponential of degree $p$, which covers the forward Euler method, every two-stage second-order method, SSPRK(3,3), and the classical fourth-order method, then
 $\eta^E_{p+2}=-(n-1)(p+1)/\big(n(n^p-1)(p+2)!\big)<0$; since $G_{p+2}(0)=1-2^{-(p+1)}>0$, at $\varpi=0$ the extrapolate is locally non-CPTP for even $p$ and locally CPTP for odd $p$.
\end{enumerate}
\end{corollary}

\begin{proof}
$E_n$ fixes $\e^z$. Write $\delta(w)=\eta_{p+1}w^{p+1}+\eta_{p+2}w^{p+2}+O(w^{p+3})$, so that $R(z/n)^n=\big(\e^{z/n}+\delta(z/n)\big)^n=\e^z+n\,\e^{(n-1)z/n}\delta(z/n)+O(z^{2p+2})$ with $2p+2\ge p+3$. Expanding $\e^{(n-1)z/n}=1+(n-1)z/n+O(z^2)$,
\[
 n\,\e^{(n-1)z/n}\delta(z/n)=\frac{\eta_{p+1}}{n^p}z^{p+1}+\frac{(n-1)\eta_{p+1}+\eta_{p+2}}{n^{p+1}}z^{p+2}+O(z^{p+3}).
\]
Multiplying by $n^p$, subtracting $R(z)-\e^z$, and dividing by $n^p-1$ cancels the $z^{p+1}$ term and leaves \eqref{eq:extrapolation_defect}. Item~2 is Theorem~\ref{thm:defect} with $m=p+2$, since $E_nR$ has real coefficients. Item~3 follows from $\eta_{p+1}=-1/(p+1)!$, $\eta_{p+2}=-1/(p+2)!$, and $(-1)^{p+2}=(-1)^p$.
\end{proof}

For $p=4$ and $n=2$ the corollary gives $\eta^E_6=-1/4320$ and, with $G_6$ from \eqref{eq:G_factors}, an extrapolate of the classical method that is locally non-CPTP for $\varpi<{(2-\sqrt{33}/4)^{1/2}}$, CPTP between the two zeros of $G_6$, and non-CPTP above ${(2+\sqrt{33}/4)^{1/2}}$; the Dormand--Prince principal formula has the opposite sequence because it has the same $G_6$ but the opposite sign of $\eta_6$. Section~\ref{sec:candidates} upgrades the local statement at $\varpi=0$ to a global one. The exponent in $E_n$ is the linear order $p$, which coincides with the classical order for the methods used here; for a method whose linear order exceeds its classical order the standard Richardson weights use the classical order, and the same computation applies with that exponent. The case $\eta_{p+1}=\eta_{p+2}$ of item~1 occurs when $R(z)\e^{-z}-1$ has no $z^{p+2}$ term, as for the symmetric implicit midpoint rule, and then extrapolation gains two orders; the per-method values of $\eta^E_{p+2}$ for $n=2,3,4$ are in Online Resource~1, Table~S2. The simplest extrapolate is that of the forward Euler method: $E_2R(z)=1+z+z^2/2$ gains one order and is the stability function of every two-stage second-order method, so its admissible set on nonrotating one-way damping is $[0,2]$ (Online Resource~1, Section~9).
\section{Exact classification for the classical fourth-order method}\label{sec:rk4}

On the boundary rays $\theta\in\{0,1\}$, $u=0$, $\nu=\pm\varpi x$ the classical fourth-order method, $R_4(z)=1+z+z^2/2+z^3/6+z^4/24$, admits a complete description of $\Adm_4(\varpi)$. Its population multiplier satisfies
\begin{equation}
 24R_4(-x)=(x^2-2x)^2+8(x^2-3x+3)>0\qquad(x\ge0),
 \label{eq:rk4_population}
\end{equation}
so $a\ge0$ never binds, while the ceiling $a\le1$ first binds at
\begin{equation}
 \alpha_4=2.785293563405\ldots,\qquad \alpha_4^3-4\alpha_4^2+12\alpha_4-24=0,
 \label{eq:alpha4}
\end{equation}
the unique positive root of that cubic. Set $q=1+4\varpi^2$ and
\begin{equation}
\begin{aligned}
 q_\pm&=\frac{123\pm11\sqrt{33}}{16},\qquad
 \varpi_\pm=\frac{\sqrt{107\pm11\sqrt{33}}}{8},\\
 \varpi_-&=0.82736225145\ldots,\qquad \varpi_+=1.63071202389\ldots,
\end{aligned}
 \label{eq:qcrit}
\end{equation}
so that $q=q_\pm$ exactly at $\varpi=\varpi_\pm$. The margin factors for every frequency, and the factorization carries the whole classification.

\begin{proposition}[Margin factorization]\label{prop:rk4factor}
On the laboratory-frame boundary rays $\theta\in\{0,1\}$, $u=0$, $\nu=\pm\varpi x$, for the classical method $R_4$ and $q=1+4\varpi^2$,
\begin{equation}
 M_4(x;\varpi)=-\frac{x^5q}{147456}\,\mathcal C_q(qx),\qquad
 \mathcal C_q(y)=y^3-16y^2-32(q-6)y+384(q-4),
 \label{eq:rk4_factor}
\end{equation}
with discriminant $\operatorname{disc}_y\mathcal C_q=16384\,(q-4)\,(8q^2-123q+348)$, which vanishes exactly at $q=4$ and $q=q_\pm$. The cubic $q\mapsto\mathcal C_q(q\alpha_4)$, whose coefficients lie in $\mathbb Q(\alpha_4)$, has exactly one real root, $q_c=3.5282058982\ldots$ ($\varpi_c=0.7950166505\ldots$), and it lies in $(1,q_-)$.
\end{proposition}

\begin{proof}
Expanding $R_4(-x)-|R_4(-\alpha x)|^2$ with $\alpha=\tfrac12-\ii\varpi$ and collecting powers of $x$ gives \eqref{eq:rk4_factor}, and the discriminant is a direct computation. The discriminant of the cubic $q\mapsto\mathcal C_q(q\alpha_4)$, reduced modulo $\alpha_4^3-4\alpha_4^2+12\alpha_4-24$, equals $-1774190592\alpha_4^2+3472883712\alpha_4+3170893824$, which is negative, the sign being decided on an isolating interval of $\alpha_4$, so the cubic has exactly one real root. Its resultant with the minimal polynomial of $\alpha_4$ is an integer polynomial of degree nine whose unique real root a Sturm sequence isolates in $(1,q_-)$, and the sign change of $\mathcal C_q(q\alpha_4)$ across that bracket identifies the root as $q_c$ (Online Resource~1, Section~3).
\end{proof}

The critical cubics factor as
\begin{equation}
\begin{aligned}
 \mathcal C_{q_-}(y)&=(y-9+\sqrt{33})^2(y-2\sqrt{33}+2),\\
 \mathcal C_4(y)&=y(y-8)^2,\\
 \mathcal C_{q_+}(y)&=(y-9-\sqrt{33})^2(y+2+2\sqrt{33}),
\end{aligned}
 \label{eq:critical_cubics}
\end{equation}
by direct expansion; these factorizations settle the equality cases of the classification.

\begin{theorem}[Classification of $\Adm_4(\varpi)$]\label{thm:rk4}
On the rays of Proposition~\ref{prop:rk4factor} write $\Adm_4(\varpi)$ for the admissible set $\Adm_{R_4}(\varpi)$ of the direct candidate at $(\theta,\kappa)\in\{(0,0),(1,0)\}$, and let $y_1\le y_2\le y_3$ denote the real roots of $\mathcal C_q$ when there are three, and $y_3$ its real root when there is one. Then:
\begin{enumerate}
 \item For $1\le q<q_-$: $\Adm_4=[0,\min\{\alpha_4,y_3/q\}]$; explicitly, $\Adm_4=[0,\alpha_4]$ for $q\le q_c$ and $\Adm_4=[0,y_3/q]$ with $y_3/q<\alpha_4$ for $q_c<q<q_-$. At $q=q_-$ the same formula holds and the double root $y=9-\sqrt{33}$ is an interior zero of the margin.
 \item For $q_-<q<4$: $\mathcal C_q$ has three positive roots and $\Adm_4=[0,y_1/q]\cup[y_2/q,y_3/q]$.
 \item At $q=4$ ($\varpi=\sqrt3/2$): $\Adm_4=\{0,2\}$.
 \item For $4<q<q_+$: $\Adm_4=\{0\}$.
 \item At $q=q_+$: $\Adm_4=\{0,(9+\sqrt{33})/q_+\}$.
 \item For $q>q_+$: $\mathcal C_q$ has one negative and two positive roots $y_2<y_3$, and $\Adm_4=\{0\}\cup[y_2/q,y_3/q]$.
\end{enumerate}
All endpoints in items 2, 5, and 6 lie below $\alpha_4$, and the isolated points in items 3 and 5 are exact members of $\Adm_4$ at which the margin vanishes.
\end{theorem}

\begin{proof}
For $x>0$, \eqref{eq:rk4_factor} makes the block condition $M_4\ge0$ equivalent to $\mathcal C_q(qx)\le0$, and \eqref{eq:rk4_population} reduces the population conditions to $x\le\alpha_4$. The sign of $\mathcal C_q$ at $y=0$ is that of $384(q-4)$. By the discriminant, $\mathcal C_q$ has one real root for $q<q_-$ and for $4<q<q_+$, and three real roots (with multiplicity) for $q_-\le q\le4$ and $q\ge q_+$. By Vieta, the roots have sum $16$, pairwise sum $32(6-q)$, and product $-384(q-4)$. For $q\le4$ no root is negative, because for $y<0$ every term of $\mathcal C_q(y)=y^3-16y^2+32(6-q)y+384(q-4)$ is nonpositive and the first is negative; hence with three real roots ($q_-\le q\le4$) all are nonnegative (strictly positive for $q<4$), and with one real root ($q<q_-$) it is positive; for $q>4$ the product is negative, so the single real root ($4<q<q_+$) is negative, and with three real roots ($q\ge q_+$) exactly one is negative. Hence for $4<q<q_+$ the cubic is positive on $y>0$ and no positive step is admissible, which is item~4; for $q>q_+$ the margin is nonnegative exactly on $[y_2/q,y_3/q]$, for $q_-<q<4$ exactly on $[0,y_1/q]\cup[y_2/q,y_3/q]$, and for $q<q_-$ exactly on $[0,y_3/q]$. Intersecting with $x\le\alpha_4$ gives items~1, 2, and 6 once the position of $y_3/q$ relative to $\alpha_4$ is known. Put $g(q)=\mathcal C_q(q\alpha_4)$; by Proposition~\ref{prop:rk4factor}, $g<0$ on $[1,q_c)$ and $g>0$ on $(q_c,\infty)$. For $q<10/3$ the derivative $\mathcal C_q'$ has no real zero, its discriminant being $128(3q-10)$, so $\mathcal C_q$ is increasing and $\operatorname{sgn}g(q)=\operatorname{sgn}(q\alpha_4-y_3)$. For $q\ge10/3$ let $y^*=\big(16+\sqrt{96q-320}\big)/3$ be the larger critical point of $\mathcal C_q$; then $q\alpha_4>y^*$ for every such $q$, because $3\alpha_4q-16\ge10\alpha_4-16>0$ and $(3\alpha_4q-16)^2-(96q-320)$ is a quadratic in $q$ with positive leading coefficient and discriminant $-2304(\alpha_4-2)(5\alpha_4+2)<0$. Since $\mathcal C_q$ is increasing on $[y^*,\infty)$, $y_3\ge y^*$ whenever there are three real roots, and $\mathcal C_q>0$ on $(y_3,\infty)$ in every case, $g(q)>0$ implies $q\alpha_4>y_3$ and $g(q)<0$ implies $q\alpha_4<y_3$. Hence $y_3/q>\alpha_4$ on $[1,q_c)$, $y_3/q=\alpha_4$ at $q_c$, and $y_3/q<\alpha_4$ on $(q_c,\infty)\supset[q_-,4]\cup[q_+,\infty)$, which gives the crossover in item~1 and places every endpoint of items 2, 5, and 6 below $\alpha_4$; the sign charts of every regime are listed in Online Resource~1 (Section~3). The equality cases follow from \eqref{eq:critical_cubics}: at $q=4$, $\mathcal C_4\le0$ on $y\ge0$ only at $y\in\{0,8\}$, giving the isolated step $x=2$; at $q=q_+$ the cubic is nonnegative on $y\ge0$ with the single zero $y=9+\sqrt{33}$; at $q=q_-$ the double root lies inside $[0,y_3]$, so the margin touches zero at an interior point without leaving the set.
\end{proof}

Figure~\ref{fig:rk4} shows two forms of reentrance: a detached interval that coexists with the attached one just below $\varpi=\sqrt3/2$, and a detached window born at $\varpi_+$ after a range with no positive admissible step. The band $\varpi_-<\varpi<\sqrt3/2$ is only $0.0386631\ldots$ wide and is easily missed by a coarse frequency grid. At $\sqrt3/2$ the components shrink to the points $0$ and $2$, and for $\sqrt3/2<\varpi<\varpi_+$ no positive step is admissible. The birth at $\varpi_+$ is a saddle-node bifurcation: the two positive roots separate as $O\big((\varpi-\varpi_+)^{1/2}\big)$ (Online Resource~1, Fig.~S2a), so companion-matrix roots are ill-conditioned nearby. This is why the classification rests on exact Sturm isolation and sign evaluation, with floating-point roots used only for plotting. The isolated steps at $q=4$ and $q=q_+$ belong to the channel set, but they have zero margin and every other step near them is non-CPTP, so a step-size controller cannot use them as targets (Section~\ref{sec:guard}).

\begin{figure}[tbp]
\centering
\includegraphics[width=\textwidth]{Fig1.eps}
\caption{Classification of $\Adm_4(\varpi)$ on the laboratory-frame one-way, zero-dephasing ray, with $x=\Gamma h$ and $\varpi=|\omega|/\Gamma$: attached and detached admissible components as functions of $\varpi$, the population ceiling $\alpha_4$, and the critical frequencies $\varpi_c$, $\varpi_-$, $\sqrt3/2$, and $\varpi_+$ of Theorem~\ref{thm:rk4}, marked by thin vertical lines ($\varpi_+$ in the main panel, the other three in the inset). The inset shows the window $0.78\le\varpi\le0.88$, $0\le x\le3$, outlined by the grey box in the main panel; it resolves the crossover at $\varpi_c$ and the reentrant band $\varpi_-<\varpi<\sqrt3/2$. Diamonds mark the isolated admissible steps at $\varpi=\sqrt3/2$ and $\varpi=\varpi_+$. All endpoints are exact roots isolated by Sturm sequences, and every endpoint belongs to the set}\label{fig:rk4}
\end{figure}

For $q\to\infty$ the two positive roots behave as $y_2\to12$ and $y_3\sim\sqrt{32q}$, hence the edges of the detached window satisfy
\begin{equation}
 x_-(\varpi)=\frac{3}{\varpi^2}+O(\varpi^{-4}),\qquad
 x_+(\varpi)=\frac{2\sqrt2}{\varpi}+O(\varpi^{-2}).
 \label{eq:window_asymptotics}
\end{equation}
To leading order the upper edge is the classical imaginary-axis absolute-stability limit of the method, because $|R_4(\ii\nu)|^2-1=\nu^6(\nu^2-8)/576$ vanishes at $\nu=2\sqrt2$; the lower edge has no scalar counterpart and is produced by the coupled population--coherence inequality. In dimensional terms, with $\nu=\omega h$, the window is $3\Gamma/\omega^2\lesssim h\lesssim2\sqrt2/|\omega|$: at high frequency the direct laboratory-frame method is a channel only for steps that resolve the precession, and then only on a window whose lower edge shrinks like the damping rate over the squared frequency.
\section{Representation and candidate construction}\label{sec:candidates}

\subsection{Rotating frame}\label{sec:frame}

Let $\tilde\rho(t)=\e^{\ii Ht}\rho(t)\e^{-\ii Ht}$. Because $\e^{\ii Ht}\sigma_\pm\e^{-\ii Ht}=\e^{\pm\ii\omega t}\sigma_\pm$ and $\mathcal D[\e^{\ii\phi}L]=\mathcal D[L]$, the transformed generator $\Ldis$ is \eqref{eq:gksl} with $\omega=0$: the interaction picture removes the coherent precession exactly and leaves the dissipators unchanged. In superoperator form $\Lgen=\Lgen_H+\Ldis$ with $[\Lgen_H,\Ldis]=0$, so $\e^{h\Lgen}=\Uh_h\,\e^{h\Ldis}$, where $\Uh_h(Y)=\e^{-\ii Hh}Y\e^{\ii Hh}$ is the unitary channel of the back-rotation. The \emph{rotating-frame candidate} is the map obtained by applying the RK method to the transformed equation and rotating back exactly,
\begin{equation}
 \Phi^{\rm RF}_R(x)=\Uh_h\circ R(h\Ldis).
 \label{eq:rf_candidate}
\end{equation}

\begin{proposition}[Representation dependence]\label{prop:frame}
(i) $\Phi^{\rm RF}_R(x)$ is CPTP if and only if $R(h\Ldis)$ is CPTP; its admissible set is therefore $\Adm_R(0)$ for every $\omega$. (ii) Let $\omega\ne0$ and let $R$ be a polynomial of degree $d\ge1$, which covers every explicit method. Then the laboratory-frame candidate $R(h\Lgen)$ and the rotating-frame candidate \eqref{eq:rf_candidate} are different linear maps: their coherence multipliers have different moduli at all but finitely many steps of any bounded interval. The same holds for the backward Euler and implicit midpoint stability functions.
\end{proposition}

\begin{proof}
(i) For every linear map $\Phi$ on $M_2(\mathbb C)$ and every unitary channel $\Uh$, $\Phi$ is completely positive if and only if $\Uh\circ\Phi$ is: compositions of completely positive maps are completely positive, and $\Uh^{-1}(Y)=\e^{\ii Hh}Y\e^{-\ii Hh}$ is itself a unitary channel, so $\Phi=\Uh^{-1}\circ(\Uh\circ\Phi)$. Since $\Uh_h$ preserves the trace and the multipliers of $R(h\Ldis)$ are the real numbers $R(-x)$ and $R(-x/2-u)$, which do not involve $\omega$, (i) follows. (ii) The coherence multiplier of \eqref{eq:rf_candidate} is $\e^{\ii\nu}R(-x/2-u)$ and that of $R(h\Lgen)$ is $R(-x/2-u+\ii\nu)$. With $u=\kappa x$ and $\nu=\varpi x$ their squared moduli are polynomials in $x$ whose coefficients of $x^{2d}$ are $r_d^2(\tfrac12+\kappa)^{2d}$ and $r_d^2\big((\tfrac12+\kappa)^2+\varpi^2\big)^d$, where $r_d\ne0$ is the leading coefficient of $R$; for $\varpi\ne0$ these differ, so the two polynomials are distinct and agree at finitely many $x$. For $R(z)=1/(1-z)$ and $R(z)=(1+z/2)/(1-z/2)$ the differences of the squared moduli are $-\nu^2/\big((1+s)^2((1+s)^2+\nu^2)\big)$ and $8\nu^2s/\big((2+s)^2((2+s)^2+\nu^2)\big)$ with $s=x/2+u$, which vanish for $\nu\ne0$ only at $s=0$.
\end{proof}

For the classical fourth-order method on the boundary rays the rotating-frame candidate is therefore admissible exactly on $[0,\alpha_4]$ at every precession frequency, whereas the laboratory-frame candidate has the frequency-dependent sets of Theorem~\ref{thm:rk4}. The detached high-frequency components, the empty regime, and the isolated steps belong to the representation as much as to the method: complete positivity is invariant under unitary conjugation, but discretizing before or after the transformation gives different maps. A static frame removes only a commuting precession; it cannot remove the noncommuting perturbation of Section~\ref{sec:noncommuting}.

\subsection{Composition and substep ladders}\label{sec:composition}

A rejected step is usually halved. On the boundary rays, when the substep population multiplier is nonnegative, substepping purely dilates the admissible set.

\begin{proposition}[Composition and dilation]\label{prop:composition}
For two maps of the form \eqref{eq:direct_map} with the same $\theta$,
\begin{equation}
 \Phi_\theta(a_2,c_2)\circ\Phi_\theta(a_1,c_1)=\Phi_\theta(a_2a_1,\,c_2c_1),
 \label{eq:composition}
\end{equation}
so the equal-substep candidate has population and coherence multipliers ${a_n=R(-x/n)^n}$ and ${c_n=R\bigl(-(x/2+u-\ii\nu)/n\bigr)^n}$. On the boundary rays $\theta\in\{0,1\}$, $u=0$, if the substep population multiplier $a_0=R(-x/n)$ is nonnegative, then $\Phi_R^{(n)}(x)$ is CPTP if and only if $\Phi_R(x/n)$ is. For the classical fourth-order method, whose population multiplier is positive by \eqref{eq:rk4_population}, $\Adm_4^{(n)}(\varpi)=n\,\Adm_4(\varpi)$ for every $n\ge1$ and every $\varpi$.
\end{proposition}

\begin{proof}
Substituting \eqref{eq:direct_map} into itself shows that the population deviation from the common fixed point $\theta$ is multiplied by $a_2a_1$ and the coherence by $c_2c_1$. If $a_0\ge0$ then $a_0^n\le1$ exactly when $a_0\le1$, and $a_0^n-|c_0|^{2n}=(a_0-|c_0|^2)\sum_{j=0}^{n-1}a_0^{n-1-j}|c_0|^{2j}$ has the sign of $a_0-|c_0|^2$ whenever the sum is positive, which fails only in the trivial case $a_0=c_0=0$, where both maps are CPTP.
\end{proof}

When $R(-x/n)\ge0$, composition therefore neither repairs a failed substep nor spoils an admissible one on these rays (the zero-precession case was observed in \cite{PierpointBernardLiu2026Note}). Without that condition it can repair one: at $\varpi=0$ two forward Euler substeps give a channel on $\{0\}\cup[8-4\sqrt2,4]$, although a single forward Euler step is a channel only at $x=0$. For the classical method substepping dilates the admissible set without changing its shape. For a detached component this has a sharp consequence for step doubling, which builds a coarse full-step candidate and a fine two-substep candidate over the same interval \cite{Shampine1985Doubling}.

\begin{proposition}[Substep ladder]\label{prop:ladder}
Suppose $\Adm_R$ has a positive-step component $\mathcal I=[x_-,x_+]$ with $0<x_-<x_+$ and that $\Adm_R^{(n)}=n\Adm_R$. Then the corresponding components $\mathcal I^{(n)}=n\mathcal I$ and $\mathcal I^{(n+1)}$ intersect if and only if $x_+/x_-\ge(n+1)/n$. In particular the full-step and two-half-step components are disjoint if and only if $x_+/x_-<2$, which holds if and only if there are initial proposals from which repeated halving skips $\mathcal I$; a fixed reduction factor $\mu\in(0,1)$ reaches $\mathcal I$ from every proposal above it if and only if $\mu\ge x_-/x_+$.
\end{proposition}

\begin{proof}
$\mathcal I^{(n)}=[nx_-,nx_+]$ and $\mathcal I^{(n+1)}={[(n+1)x_-,(n+1)x_+]}$ intersect exactly when ${(n+1)x_-}\le nx_+$. Consecutive proposals $\mu^jx_0>x_+$ and $\mu^{j+1}x_0<x_-$ from some initial proposal $x_0$ exist exactly when $\mu<x_-/x_+$.
\end{proof}

\begin{example}[Disjoint step-doubling candidates at $\varpi=2$]\label{ex:varpi2}
Exact isolation gives
\begin{equation}
\begin{aligned}
 \Adm_4(2)&=\{0\}\cup[0.7447824278959\ldots,\,1.2703346269738\ldots],\\
 \Adm_4^{(2)}(2)&=\{0\}\cup[1.4895648557918\ldots,\,2.5406692539477\ldots],
\end{aligned}
 \label{eq:varpi2_sets}
\end{equation}
whose positive-step components are disjoint because $x_+/x_-=1.7056\ldots<2$: at total step $x=1$ the full-step candidate is a channel and the two-half-step candidate is not, and at $x=2$ the statuses reverse (Fig.~\ref{fig:candidates}a). The pair $n=1,2$ is the most restrictive rung of the ladder; at this frequency $\mathcal I^{(n)}$ and $\mathcal I^{(n+1)}$ intersect for every $n\ge2$. For the high-frequency window, disjointness and the halving-skip condition end together at $\varpi_{1/2}=2.248254604514980\ldots$, where $x_+/x_-=2$. The condition that one root of $\mathcal C_q$ be twice another is detected by the resultant $\operatorname{Res}_y(\mathcal C_q(y),\mathcal C_q(2y))=6291456\,(q-4)\,(72q^3-2071q^2+12552q-21744)$, and $\varpi_{1/2}$, where $y_3=2y_2$, corresponds to the unique root of the cubic factor above $q_+$. Its two other roots lie in the reentrant band: at $q=3.7645\ldots$ ($\varpi=0.8313\ldots$) the relation is $y_2=2y_1$, between the attached and detached components, whereas at $q=3.7807\ldots$ ($\varpi=0.8337\ldots$) the detached component itself has $x_+/x_-=2$, so for $3.7807\ldots<q<4$ its full-step and two-half-step copies are also disjoint.
\end{example}

\begin{figure}[tbp]
\centering
\includegraphics[width=\textwidth]{Fig2.eps}
\caption{Maps built over the same interval by the classical fourth-order method. \textbf{a}~At $\varpi=2$ the positive-step admissible components of the full-step candidate (one step of size $x$) and of the two-half-step candidate are disjoint. Both components are detached (Theorem~\ref{thm:rk4}, item~6), and both sets also contain $x=0$. At the marked total steps $x=1$ and $x=2$ the two candidates exchange roles as the channel. \textbf{b}~On nonrotating one-way damping the coarse and fine candidates are channels up to $x=\alpha_4$, whereas the Richardson extrapolate has a strictly negative Choi margin throughout that interval (Theorem~\ref{thm:richardson}). The curve in b is evaluated from the exact polynomial}\label{fig:candidates}
\end{figure}

Many step-doubling codes advance the fine candidate and use the coarse one only to estimate the error. Only the advanced candidate needs a channel certificate, but a certificate for one does not transfer to the other, and requiring both to be channels leaves no admissible positive step at $\varpi=2$.

\subsection{Richardson extrapolation}\label{sec:richardson}

The set of CPTP maps is convex, so a convex combination of channels is a channel; trace preservation survives every affine combination whose weights sum to one, but a negative weight can leave the channel set \cite{Choi1975CP}. For the classical fourth-order method let $\Phi_{\rm c}(x)=\Phi_4(x)$ and $\Phi_{\rm f}(x)=\Phi_4(x/2)^2$ be the coarse and fine candidates and
\begin{equation}
 \Phi_{\rm ext}(x)=\frac{16\,\Phi_{\rm f}(x)-\Phi_{\rm c}(x)}{15}
 \label{eq:richardson_map}
\end{equation}
the standard Richardson extrapolate ($n=2$, $p=4$ in Definition~\ref{def:candidate}). By Proposition~\ref{prop:composition} and the affine structure of \eqref{eq:direct_map}, $\Phi_{\rm ext}$ is the phase-covariant map generated by
\begin{equation}
\begin{aligned}
 R_{\rm ext}(z)&=\frac{16R_4(z/2)^2-R_4(z)}{15}=\sum_{j=0}^{5}\frac{z^j}{j!}+\frac{z^6}{864}+\frac{z^7}{8640}+\frac{z^8}{138240},\\
 R_{\rm ext}(z)-\e^z&=-\frac{z^6}{4320}+O(z^7),
\end{aligned}
 \label{eq:richardson_stability}
\end{equation}
in agreement with Corollary~\ref{cor:extrapolation}. The corollary's local statement at $\varpi=0$ becomes global: the extrapolate is a channel at no positive step where both constituents are.

\begin{theorem}[Loss of complete positivity under extrapolation]\label{thm:richardson}
On nonrotating one-way damping ($\varpi=0$, $u=0$, $\theta\in\{0,1\}$), $\Phi_{\rm c}(x)$ is CPTP for $0\le x\le\alpha_4$ and $\Phi_{\rm f}(x)$ for $0\le x\le2\alpha_4$, whereas $\Phi_{\rm ext}(x)$ is not CPTP for any $0<x\le\alpha_4$. Its margin factors as
\begin{equation}
 M_{\rm ext}(x)=-\frac{x^6\,P_{10}(x)}{1252412463513600}=-\frac{31}{138240}\,x^6+O(x^7),
 \label{eq:richardson_margin}
\end{equation}
where $P_{10}$ is the degree-ten integer polynomial listed in Online Resource~1, whose smallest positive root $6.034523958649\ldots$ exceeds $\alpha_4$. The extrapolated candidate re-enters the channel set only on the remote interval $[6.034523958649\ldots,\,6.459127767825\ldots]$, whose upper end is the population-ceiling exit $R_{\rm ext}(-x)=1$.
\end{theorem}

\begin{proof}
$\Adm_4(0)=[0,\alpha_4]$ by Theorem~\ref{thm:rk4} and $\Adm_4^{(2)}(0)=[0,2\alpha_4]$ by Proposition~\ref{prop:composition}. Corollary~\ref{cor:extrapolation} with $p=4$, $n=2$ gives the leading coefficient $-\tfrac1{4320}G_6(0)=-\tfrac{31}{138240}$; exact substitution into $R_{\rm ext}(-x)-R_{\rm ext}(-x/2)^2$ produces the factor $P_{10}$, and a Sturm sequence places its first positive zero above $\alpha_4$, so the margin is negative on $(0,\alpha_4]$. The complete sign chart of the margin, of $R_{\rm ext}(-x)$, and of $1-R_{\rm ext}(-x)$ in Online Resource~1 isolates the remote component.
\end{proof}

Codes that advance $\Phi_{\rm f}$ and use the extrapolate only as an error indicator are unaffected. A code that advances the extrapolate must test its channel condition separately, whatever its order and however physical its constituents. At $\omega=0$ the two frames give the same map, so the failure occurs in both; the frequency-dependent orientation bands of Section~\ref{sec:extrapolation} concern the laboratory-frame candidate. Figure~\ref{fig:candidates}b shows the scaled extrapolation margin, negative throughout $(0,\alpha_4]$.

\begin{remark}[Convergence from outside the channel set]\label{rem:fixed_horizon}
None of this conflicts with convergence. On nonrotating one-way damping take $N$ equal steps $x_N=X/N$ to a fixed horizon $X$ and let $R(z)-\e^z=\eta_mz^m+O(z^{m+1})$. Then the margin of the composed map satisfies
$M_N(X)=K_m\,X\e^{-X}(X/N)^{m-1}+O(N^{-m})$ with $K_m=(-1)^m(1-2^{1-m})\eta_m$ (Online Resource~1, Section~8). When $K_m<0$ the composed candidate converges to the exact channel while remaining outside the channel set at every sufficiently fine resolution, with a negative Choi eigenvalue that decays at the global order $m-1$; forward Euler and SSPRK(3,3) are examples with exact negative margins throughout their positive-population range. Convergence of the maps is therefore compatible with non-physical maps at every finite resolution \cite{PierpointBernardLiu2026Note}.
\end{remark}
\section{Persistence and certified acceptance}\label{sec:guard}

The disconnected sets of Theorem~\ref{thm:rk4} live on a boundary ray. This section proves that they survive small buffers \eqref{eq:buffers}, certifies them at three strengths of a noncommuting perturbation, and turns the criterion into an adaptive acceptance test.

\subsection{Persistence under buffers}\label{sec:persistence}

\begin{theorem}[Relative-open persistence of disconnectedness]\label{thm:persistence}
Fix a method of order at least two, a ratio $\varpi$, and a compact step interval $\mathcal K$ contained in the stage-evaluation domain. Suppose that at $(\theta,\kappa)=(0,0)$: (i) the Choi margin has a detached positive component $[r_2,r_3]\subset\mathcal K$ bounded by simple roots; (ii) there are points $s_-\in(0,r_2)$, $s_+\in(r_2,r_3)$, and $s_{\rm out}>r_3$ in $\mathcal K$ at which the Choi margin has the signs $-,+,-$; (iii) the constraints of \eqref{eq:criterion} other than $\Fmargin\ge0$ hold with uniform strict slack on a closed neighbourhood of $[r_2,r_3]$ contained in $\mathcal K$. Then there is a relative-open neighbourhood $N$ of $(0,0)$ in the quadrant $\theta,\kappa\ge0$ such that for every $(\theta,\kappa)\in N$ with $\theta+\kappa>0$ the admissible set has at least two positive-step components, an attached interval and a detached interval near $[r_2,r_3]$. The theorem does not describe the steps between the attached interval and $s_-$.
\end{theorem}

\begin{proof}
The three strict signs, the uniform slack of the other constraints, and (for implicit methods) the open stage-evaluation condition persist on a relative-open parameter neighbourhood by continuity and compactness; the implicit function theorem continues both simple roots, the buffers \eqref{eq:buffers} supply the attached interval, and the margin stays positive on $[r_2+\varepsilon,r_3-\varepsilon]$, giving the detached component. Details are in Online Resource~1, Section~4.
\end{proof}

\begin{figure}[tbp]
\centering
\includegraphics[width=\textwidth]{Fig3.eps}
\caption{Disconnected admissible sets of the classical fourth-order method at $\varpi=2$, with $x=\Gamma h$. \textbf{a}~Exact components for five buffered phase-covariant cases $(\theta,\kappa)$; endpoints are simple roots isolated by Sturm counts. \textbf{b}~Upper endpoint of the attached component and both endpoints of the detached component for the transverse-field generator \eqref{eq:transverse} with $\theta=\kappa=10^{-3}$ at $\Omega_x/\Gamma\in\{0,0.05,0.10,0.20\}$; the markers are not joined because each coupling is certified separately}\label{fig:robustness}
\end{figure}

For the classical fourth-order method at $\varpi=2$ the hypotheses hold with $s_-=1/2$, $s_+=1$, $s_{\rm out}=3/2$. Figure~\ref{fig:robustness}a shows the exact components $[0,r_1]\cup[r_2,r_3]$ for five buffered cases (endpoints and sign samples in Online Resource~1, Section~4). The attached interval grows with the buffer as \eqref{eq:buffers} predicts, from $r_1=0.1222\ldots$ at $(\theta,\kappa)=(10^{-3},0)$ to $r_1=0.5193\ldots$ at $(0,10^{-2})$. Repeated reduction of a step $x_0$ above the attached endpoint $x_{\rm att}=r_1$ by a factor $\mu\in(0,1)$ reaches the attached interval after about $\lceil\log(x_{\rm att}/x_0)/\log\mu\rceil$ reductions, and the joint-limit estimates $x_{\rm att}\sim(48\kappa)^{1/4}/\varpi$ and $x_{\rm att}\sim(24\theta(1-\theta)/\varpi^4)^{1/3}$, from \eqref{eq:buffers} against the large-frequency defect $-\varpi^4x^5/24$, show that this number diverges as the buffer vanishes.

\subsection{A noncommuting perturbation}\label{sec:noncommuting}

The two-mode reduction relies on phase covariance; disconnectedness does not. Consider
\begin{equation}
 \Lgen_{\Omega_x}(\rho)=\Lgen(\rho)-\ii\big[\Omega_x\sigma_x/2,\rho\big],
 \label{eq:transverse}
\end{equation}
with the dissipators of \eqref{eq:gksl}. For $\Omega_x\ne0$ phase covariance is broken, no static frame removes the precession, and the Choi matrix of $R_4(h\Lgen_{\Omega_x})$ must be assembled from the full $4\times4$ superoperator (Online Resource~1, Section~4). A sign pattern in the step is then enough to force two components.

\begin{lemma}[Persistence in a parameterized family]\label{prop:noncommuting}
Let $\Lgen_\eta$ be a continuous finite-dimensional family of GKSL generators. If four steps $0<h_a<h_g<h_d<h_o$ give strict signs $+,-,+,-$ of the smallest Choi eigenvalue of $R_4(h\Lgen_{\eta_0})$, then the same sign pattern, and hence at least two positive-step completely positive components, persists for all $\eta$ in a neighbourhood of $\eta_0$.
\end{lemma}

\begin{proof}
The Choi matrix of $R_4(h\Lgen_\eta)$ is Hermitian and depends continuously on $(h,\eta)$, and so does its smallest eigenvalue; hence four strict signs persist on a common neighbourhood, and two positive samples separated by a negative one cannot lie in one component.
\end{proof}

For the generator \eqref{eq:transverse} with $\Gamma=1$, $\theta=\kappa=10^{-3}$, $\omega/\Gamma=2$, and $\Omega_x/\Gamma=0.1$, the smallest Choi eigenvalues obtained from exact characteristic-polynomial brackets are
\mbox{$9.5156\times10^{-5}$}, \mbox{$-2.7778\times10^{-3}$}, \mbox{$5.8959\times10^{-4}$}, and \mbox{$-3.9227\times10^{-1}$} at $h=0.1$, $0.5$, $0.9$, and $1.4$, and exact isolation of the roots of every principal minor yields
\begin{equation}
\begin{aligned}
 \{h\in[0,1.5]:\ R_4(h\Lgen_{0.1})\ \text{is CPTP}\}
 ={}&[0,\,0.26357690\ldots]\\
 &\cup[0.73588381\ldots,\,1.26847247\ldots].
\end{aligned}
 \label{eq:noncommuting_set}
\end{equation}
Repeating the exact computation at $\Omega_x/\Gamma\in\{0,0.05,0.10,0.20\}$ gives an attached and a detached component in each case (Fig.~\ref{fig:robustness}b; endpoints in Online Resource~1). These four certificates, three of them at nonzero field, are pointwise; Lemma~\ref{prop:noncommuting} extends each to a neighbourhood of its coupling.

\subsection{Certified acceptance of the executed step}\label{sec:certified}

Adaptive step-size control proposes steps from local error estimates \cite{DormandPrince1980,HairerNorsettWanner1993,Shampine1985Doubling,Soderlind2002}. An error estimate cannot see whether the update is a channel, so we add a second test: a complete-positivity certificate for the map about to be advanced. A failed certificate rejects that map without feeding an artificial error ratio into the controller.

\paragraph{The certified map.} The certificate concerns the linear map $\Phi_{\rm spec}(x)$ that the construction rule of Definition~\ref{def:candidate} assigns to the exact rational value of the executed double-precision step. The test has three outcomes: \PASS\ and \FAIL\ are exact decisions of \eqref{eq:criterion} and of the stage-evaluation condition, reached by rational arithmetic and, where roots must be located, by Sturm isolation; \UNCERTAIN\ is returned whenever a sign, a root identity, an ordering, or a root separation cannot be settled at the working precision. An undecided step is never accepted. The double-precision matrix multiplied into the accumulated propagator is a rounding of $\Phi_{\rm spec}(x)$, and the certificate does not enclose the floating-point products that follow.

Exact arithmetic is forced by the geometry near the boundary. On the one-way ray the margin of the classical method is $-x^5q\,\mathcal C_q(qx)/147456$ by \eqref{eq:rk4_factor}, a fifth-order quantity that floating-point evaluation of $R_4(-x)-|R_4(-\alpha x)|^2$ loses to cancellation as $x\to0$ (Online Resource~1, Fig.~S2c); a tolerance-based eigenvalue test then accepts non-channels at small steps. Near an endpoint, decimal-to-binary conversion alone can flip the sign: at $\varpi=2$ the decimal $x=1.270334626973889$ has margin \mbox{$+1.43\times10^{-16}$}, whereas the double-precision number it denotes has margin \mbox{$-9.45\times10^{-17}$}.

\paragraph{Projection.} The natural replacement for a rejected step is a smaller step of the same construction. Let $\Adm^{\rm cand}(x_{\rm prop})$ be the admissible set of the proposed construction on $(0,x_{\rm prop}]$. On a phase-covariant ray the constraints \eqref{eq:criterion} of the polynomial candidates in Definition~\ref{def:candidate} are polynomial inequalities with rational coefficients, so the components of $\Adm^{\rm cand}(x_{\rm prop})$ are located by exact root isolation (Online Resource~1, Section~7). A component of positive width contains rational steps but need not contain a double-precision number that a finite search finds. A projection is therefore accepted only at a double-precision step that lies in the component and at which the candidate has been certified again. Otherwise the outcome is \UNCERTAIN, and when no component of positive width lies below the proposal the step needs a different construction. Admissibility does not imply accuracy, so the error is recomputed at the executed step.

\begin{algorithm}[t]
\caption{Adaptive step with exact complete-positivity acceptance}\label{alg:guard}
\begin{algorithmic}[1]
\Require generator ray, RK method, candidate construction rule, step proposal $x_{\rm prop}$ from the error controller, tolerance, fallback constructions (other representations, or a propagator that is a channel by construction)
\Loop
  \State construct the candidate to be advanced at $x_{\rm prop}$ and its local-error estimate
  \If{the error estimate exceeds the tolerance}
    \State update $x_{\rm prop}$ from the error estimate
  \ElsIf{the exact certificate of the candidate at $x_{\rm prop}$ returns \PASS}
    \State advance this candidate, pass its error to the controller, and \Return
  \ElsIf{\FAIL\ and $\Adm^{\rm cand}(x_{\rm prop})$ has a positive-width component}
    \State project $x_{\rm prop}$ to a certified double-precision step inside it
  \ElsIf{\FAIL}
    \State switch to the next fallback construction
  \ElsIf{\UNCERTAIN}
    \State raise the working precision, or switch to the next fallback construction
  \EndIf
\EndLoop
\end{algorithmic}
\end{algorithm}

Algorithm~\ref{alg:guard} keeps error control and complete-positivity control separate: the certificate ignores the error estimate, a rejection on complete-positivity grounds leaves the controller's error history unchanged, the error is recomputed at any projected candidate before acceptance, and a change of representation or propagator discards an error ratio that belongs to a different candidate family. In items 3--6 of Theorem~\ref{thm:rk4} there is no attached interval, so on the boundary rays a proposal below the first positive admissible step (if any) has no smaller admissible direct step and the algorithm must switch; with any buffer an attached interval returns by \eqref{eq:buffers}, and projection becomes a cost decision. A fallback is accepted on the strength of its construction only when it is the exact exponential or a composition of exact subflows with a nonnegative time argument, and only after the matrix realized in double precision passes a residual check against that construction. Any other fallback, such as the rotating-frame candidate below, is certified in the same way on the next pass. The loop ends when a candidate passes both tests, and a failed projection search counts as \UNCERTAIN. If the error test passes at sufficiently small steps, the loop terminates once the fallback is a channel by construction, whose error estimate is zero, or has an attached admissible interval, as the rotating-frame candidate has on the boundary rays by Proposition~\ref{prop:frame}: every later rejection then leaves a positive-width component below the proposal.

\subsection{Adaptive tests}\label{sec:adaptive_tests}

The experiments test whether every accepted update is a channel. Step doubling estimates the error from the coarse and fine candidates of the classical fourth-order method and advances the fine candidate; the error measure is the normalized Choi trace distance $\tfrac1{15}\cdot\tfrac12\big\|\big(J[\Phi_{\rm c}]-J[\Phi_{\rm f}]\big)/2\big\|_1$, the proposal rule is a proportional--integral controller with a safety factor of $0.9$, and the exact map $\e^{X\Lgen}$ over the whole horizon is the reference for the global error. Four policies are compared on the boundary case $\theta=\kappa=0$, $\varpi=2$, horizon $X=2$, initial step $1.4$: error-only laboratory-frame stepping, Algorithm~\ref{alg:guard} with the rotating-frame candidate as fallback (the guard in Table~\ref{tab:adaptive}), rotating-frame stepping throughout, and Strang splitting~\cite{Strang1968} of the exact coherent and dissipative subflows, which is exact for this commuting generator and serves as a roundoff-floor reference. Buffered ($\theta=0$, $\kappa=10^{-3}$, $\varpi=2$, initial step $0.3$) and near-saddle ($\theta=0$, $\kappa=10^{-3}$, $\varpi=1.65$, initial step $1.0$, horizon $X=1$) cases show the guard keeping direct steps and projecting rejected ones. Ray parameters are entered as double-precision numbers and each certificate refers to their exact rational values. Table~\ref{tab:adaptive} reports accuracy and physicality separately; the accuracy--work relation over the tolerance sweep and the accepted candidate types are plotted in Online Resource~1, Fig.~S4.

\begin{table}[tbp]
\caption{Adaptive tests. Choi error: global normalized Choi trace distance to the exact map. Min.\ eig.: smallest double-precision Choi eigenvalue of the executed matrices over all accepted steps (acceptance itself was decided exactly); $\ge-10^{-15}$ is roundoff level, and a more negative value means that a non-channel was accepted. d/p/f: accepted candidates of guarded runs by type, direct/projected/fallback; -- for policies without the guard. RK stages: right-hand-side evaluations of the advanced and estimator candidates; the Strang row uses 18 exponential actions instead. Tolerance $10^{-3}$ in every row except buffered (2), which uses \mbox{$3\times10^{-4}$}; the near-saddle row has horizon $X=1$}\label{tab:adaptive}
\footnotesize
\begin{tabular*}{\textwidth}{@{\extracolsep\fill}llcccr}
\toprule
Case & Policy & Choi error & Min.\ eig. & d/p/f & RK stages \\
\midrule
boundary & error only & $7.92\times10^{-4}$ & $-1.0\times10^{-3}$ & -- & 72 \\
boundary & guard, rotating fallback & $7.45\times10^{-6}$ & $\ge-10^{-15}$ & 0/0/4 & 132 \\
boundary & rotating frame & $1.27\times10^{-4}$ & $\ge-10^{-15}$ & -- & 36 \\
boundary & Strang, exact subflows & $8.33\times10^{-17}$ & $\ge-10^{-15}$ & -- & 0 \\
buffered (1) & guard & $2.85\times10^{-4}$ & $\ge-10^{-15}$ & 3/2/0 & 84 \\
buffered (2) & guard & $2.47\times10^{-4}$ & $\ge-10^{-15}$ & 5/0/0 & 60 \\
near saddle & guard & $2.10\times10^{-4}$ & $\ge-10^{-15}$ & 1/1/0 & 48 \\
\botrule
\end{tabular*}
\end{table}

Error-only step doubling meets its tolerance on the boundary case while accepting non-channel candidates (minimum Choi eigenvalue $-10^{-3}$), as Theorem~\ref{thm:rk4} predicts: at $\varpi=2$ the fine candidate is a channel only on $[1.4895\ldots,2.5406\ldots]$, above every proposed step. The guard accepts no non-channel candidate. On the boundary case none of its rejected proposals had a positive-width component below it, so every accepted step was a rotating-frame fallback (four in Table~\ref{tab:adaptive}, fifteen over the four tolerances of Online Resource~1, Table~S3), at 1.4 to 1.8 times the stage count of error-only stepping at the same tolerance; the cost of the certificates themselves is reported separately in Online Resource~1. With a buffer the guard retains direct candidates and projects to certified interior steps (two projections in the first buffered row, one near the saddle) at global errors below $10^{-3}$. The lower stage count of the rotating-frame policy and the exactness of the Strang reference both rest on the commuting dynamics of this benchmark and give no general ranking of the methods.

\begin{remark}[Scope of the adaptive certificates]\label{rem:scope}
Each accepted step is specified by a linear map certified to be CPTP at its exact rational step. When the step sequence depends on one particular input state, however, the adaptive input--output procedure is nonlinear in that state and is not itself a reusable channel, which would require a state-independent error criterion.
\end{remark}
\section{Discussion and limitations}\label{sec:discussion}

\paragraph{Relation to classical stability notions.} On one-way damping both scalar multipliers can be contractive while $|c|^2\le a$ fails, and the admissible set can be disconnected (Theorem~\ref{thm:rk4}), whereas threshold-type positivity and contractivity results \cite{BolleyCrouzeix1978,Kraaijevanger1991Contractivity,Horvath2005Threshold,GottliebShuTadmor2001} describe the maximal attached interval. On the rays of Theorem~\ref{thm:rk4} the per-mode absolute-stability sets are intervals; the coupled inequality alone produces the disconnection. Adaptive positivity by weight perturbation \cite{NussleinRanochaKetcheson2021} and the positivity barrier for exponential RK methods \cite{OstermannVanDaele2007} constrain the solution of an ODE; here the unchanged RK map is constrained under tensor extension.

\paragraph{Relation to structure-preserving propagators.} The structure-preserving propagators cited in Section~\ref{sec:intro} \cite{RieschPiklJirauschek2020,CaoLu2025,AppeloCheng2025,HuAppeloCheng2025NPI,HuAppeloCheng2026Gregory,DelMastroAppeloCheng2026TT,RobinRouchonSellem2025,ChenEtAl2026} modify the update so that it preserves structure by construction, with the different guarantees noted there. We ask instead whether the unchanged RK update, formed directly or by substepping, extrapolation, or a change of representation, is a channel, and if not, whether a smaller certified step exists or a different representation or propagator is needed. Methods that preserve complete positivity by construction remain the natural choice when their order and cost suit the problem; the present results explain the failures that local error control cannot detect.

\paragraph{Cost.} One membership decision costs two scalar evaluations of the stability function and four inequalities, plus the stage-domain check for implicit methods; locating every component or proving emptiness is real-root isolation on polynomials of moderate degree \cite{BasuPollackRoy2006}. The two-mode reduction does not extend at that cost to a black-box or higher-dimensional Liouvillian, where a map-level Choi test is substantially larger, as the transverse-field certificates of Section~\ref{sec:noncommuting} show already for one qubit.

\paragraph{Limitations.} The theorems concern the autonomous phase-covariant qubit family, and the six-regime classification concerns the classical fourth-order method on the laboratory-frame boundary rays; other methods have their own admissible sets (Table~\ref{tab:methods}; Online Resource~1, Section~9) and reversal frequencies, to which the criterion applies unchanged. Theorem~\ref{thm:persistence} and Lemma~\ref{prop:noncommuting} are local statements, and the buffered and transverse-field sets are certified only at the listed parameters, the latter by four pointwise certificates, three at nonzero field. The disjointness at $\varpi=2$ concerns one full step against two half steps, and Theorem~\ref{thm:richardson} concerns advancing the extrapolated map. A certificate refers to the specified linear map at exact double-precision inputs, not to the accumulated floating-point product (Section~\ref{sec:certified}), and an adaptive procedure whose steps depend on the input state is not itself a channel (Remark~\ref{rem:scope}). A step that is a channel can have a larger local error than one that is not, which is why the adaptive tests report accuracy and physicality separately.
\section{Conclusion}\label{sec:conclusion}

For the autonomous phase-covariant qubit generator, the complete positivity of a Runge\nobreakdash--Kutta step is decided exactly by two values of the stability function through the known channel conditions. On the boundary rays the first defect from the exponential gives the leading margin coefficient for direct steps and Richardson extrapolates, and whenever that coefficient is nonzero its sign decides whether sufficiently small steps are channels. For the classical fourth-order method on these rays the classification is complete: a reentrant band, a regime without positive admissible steps, isolated admissible steps, and a detached high-frequency interval with edges $3/\varpi^2$ and $2\sqrt2/\varpi$ to leading order, separated by critical frequencies that are explicit algebraic numbers.

The practical rule is to certify the update that is advanced, in the representation in which it is advanced, at the step that is executed. The frame, substepping, and extrapolation each change the admissible set, and local error control detects none of these changes. For phase-covariant candidates the exact acceptance test decides membership with two evaluations of the stability function in rational arithmetic and, in the boundary tests, rejects every non-channel update that error control accepts. Non-phase-covariant and time-dependent generators are the natural next cases.

\backmatter

\bmhead{Supplementary information}
Online Resource 1 (\texttt{ESM\_1.pdf}) contains the sign charts and root certificates for Theorem~\ref{thm:rk4}, the per-method table for Corollary~\ref{cor:extrapolation}, the transverse-field construction and certificates, projection and step-control details, the adaptive tolerance sweep with its accuracy--work figure, and supporting derivations.

\bmhead{Acknowledgements}
This research was supported by the Research Computing clusters at Old Dominion University. The Wahab cluster is supported in part by National Science Foundation grant CNS-1828593, ``MRI Acquisition: A Reconfigurable Computing Infrastructure Enabling Interdisciplinary and Collaborative Research in Hampton Roads.''

\phantomsection\addcontentsline{toc}{section}{Declarations}
\section*{Declarations}

\bmhead{Funding}
This research received no specific grant from funding agencies in the public, commercial, or not-for-profit sectors.

\bmhead{Competing interests}
The authors have no competing interests to declare that are relevant to the content of this article.

\bmhead{Ethics approval and consent to participate}
Not applicable.

\bmhead{Consent for publication}
Not applicable.

\bmhead{Data availability}
All data generated or analysed for this article and Online Resource~1 are included in the code and data release cited under Code availability.

\bmhead{Code availability}
{\raggedright The code and data are available under the MIT License in release~v4.1.1 of the repository \url{https://github.com/Weyl-Center-for-Mathematical-Physics/rk-choi-complete-positivity}, archived on Zenodo at \url{https://doi.org/10.5281/zenodo.21522057}. It computes every certificate, sign chart, table, and figure in the article and Online Resource~1, and explains how to reproduce them. The earlier zero-precession amplitude-damping release is archived at \cite{PierpointBernardLiu2026Zenodo}.\par}

\bmhead{Author contributions}
Conceptualization: G.~Blake Pierpoint, Olivier Bernard; Methodology: G.~Blake Pierpoint, Olivier Bernard, Yichen Liu; Formal analysis and investigation: G.~Blake Pierpoint, Olivier Bernard, Yichen Liu; Software: G.~Blake Pierpoint, Olivier Bernard, Yichen Liu; Validation: Olivier Bernard, Yichen Liu; Data curation and visualization: G.~Blake Pierpoint; Project administration and supervision: G.~Blake Pierpoint; Writing -- original draft: G.~Blake Pierpoint; Writing -- review and editing: G.~Blake Pierpoint, Olivier Bernard, Yichen Liu. All authors approved the final manuscript.

\bmhead{Use of generative artificial intelligence}
Anthropic Claude (Claude Code) and OpenAI tools (Codex) were used for drafting and editing text, literature searches, mathematical analysis, and developing and checking code for symbolic verification, numerical experiments, figures, and packaging. The authors approved the final manuscript and take responsibility for its claims, references, and computations.

% Keep each reference intact rather than stranding a DOI continuation on the next page.
\renewcommand{\bibpreamble}{\interlinepenalty=10000}
\phantomsection\addcontentsline{toc}{section}{References}
\bibliography{references}

\end{document}
